\documentclass[11pt]{article}

\usepackage{mathtools}


\title{Actuation Matrices}
\author{W. Muneo Okiishi}
\date{2026-10-15}

\begin{document}

\maketitle

The shape change of a truss robot can be described by either
$\dot{\mathbf{q}}$, the joint velocities, or $\dot{\boldsymbol{\ell}}$,
the rate of change of the lengths of its members.

The relationship between these can be easily encoded in a matrix $\mathbf{A}$
such that

\begin{equation}
    \mathbf{A} \dot{\mathbf{q}} = \dot{\boldsymbol{\ell}}
\end{equation}

Generally, any $\dot{\mathbf{q}}$ within limits is achievable, but some
combinations of $\dot{\boldsymbol{\ell}}$ are unreachable by the physical
robot. Achievable values of $\dot{\boldsymbol{\ell}}$ must be within
$\mathcal{C}(\mathbf{A})$, the column space of $\mathbf{A}$.

For example, if $\mathbf{A}^+$ is the using the Moore--Penrose inverse
(pseudoinverse) of $\mathbf{A}$, then

\begin{equation}
    \mathbf{A} \mathbf{A}^+ \dot{\boldsymbol{\ell}} = \dot{\boldsymbol{\ell}}
    \text{ for all } \dot{\boldsymbol{\ell}} \in \mathcal{C}(\mathbf{A})
\end{equation}

Thus, only when $\dot{\boldsymbol{\ell}} \in \mathcal{C}(\mathbf{A})$ will a
$\dot{\mathbf{q}}$ constructed as $\dot{\mathbf{q}} = \mathbf{A}^+ \dot{\boldsymbol{\ell}}$
produce the original value of $\dot{\boldsymbol{\ell}}$.

In order to to evaluate whether a given $\dot{\boldsymbol{\ell}}$, we thus
want to find a matrix $\mathbf{T}$ such that

\begin{equation}
    \mathbf{T} \dot{\boldsymbol{\ell}} = \mathbf{0}
    \iff \dot{\boldsymbol{\ell}} \in \mathcal{C}(\mathbf{A})
\end{equation}

This is the case when

\begin{equation}
    \mathcal{N}(\mathbf{T}) = \mathcal{C}(\mathbf{A})
\end{equation}


\end{document}
